\documentclass[12pt,a4paper]{article}

\usepackage[utf8]{inputenc}
\usepackage[T1]{fontenc}
\usepackage[ngerman]{babel}
\usepackage{amsmath,amssymb,amsthm,mathtools}
\usepackage[colorlinks=true,linkcolor=blue,citecolor=red,urlcolor=blue]{hyperref}
\renewcommand*{\HyperDestNameFilter}[1]{\jobname-#1}
\usepackage{geometry}
\geometry{margin=2.5cm}
\emergencystretch=3em
\usepackage{enumitem}
\usepackage{booktabs}
\usepackage{graphicx}
\usepackage{needspace}
\usepackage[capitalise,noabbrev,ngerman]{cleveref}

\newtheorem{theorem}{Theorem}[section]
\newtheorem{lemma}[theorem]{Lemma}
\newtheorem{proposition}[theorem]{Proposition}
\newtheorem{corollary}[theorem]{Korollar}
\newtheorem{conjecture}[theorem]{Vermutung}
\theoremstyle{definition}
\newtheorem{definition}[theorem]{Definition}
\theoremstyle{remark}
\newtheorem{remark}[theorem]{Bemerkung}

\title{Die Riemannsche Vermutung: Conclusio\\[4pt]
\large Teil~V einer fünfteiligen Serie}
\author{Lukas Geiger\\
\small Bernau, Deutschland\\
\small ORCID: \href{https://orcid.org/0009-0005-7296-1534}{0009-0005-7296-1534}}
\date{\today}

\begin{document}
\overfullrule=0pt
\maketitle

\begin{abstract}\footnotesize
Dieses Paper schließt eine fünfteilige Untersuchung der Riemannschen
Vermutung durch Spektralanalyse, arithmetische Thermodynamik und die
Weil'sche quadratische Form ab:
\begin{description}[font=\bfseries,leftmargin=2em]
  \item[Teil~I \cite{Geiger2026partI}:] Grundlagen und Obstruktionen.
  \item[Teil~II \cite{Geiger2026partII}:] Even Dominance (Zweig~A).
  \item[Teil~III \cite{Geiger2026partIII}:] Spektrale Routen
    (Zweige~B und~Z).
  \item[Teil~IV \cite{Geiger2026partIV}:] Routenübergreifende Synthese.
  \item[Teil~V (dieses Paper):] Conclusio.
\end{description}

\noindent Wir präsentieren den kondensierten Atlas: drei unabhängige
bedingte Reduktionen der RH (Even Dominance, Hadamard Strip, CCM
Microcluster), drei strukturelle Identifikationen
($c_\Xi \approx 0.0231$, $C^* = 64\pi^2/3$ als $\lambda=5,7$-
Boundary-Moment-Benchmark mit offenem $\lambda=11$-Diskriminator,
$g_0(m) \approx 2\pi/\log\gamma_m$), zwei
bedingte Theoreme (B-1, GF1) und zwei
isolierte offene Beweis-Slots. Der Serientitel bleibt \glqq From
Landscape to Atlas\grqq{}; die Heraufstufung zu \glqq From Atlas to
Proof\grqq{} wird durch ein Fünf-Tore-Audit gesteuert, das nur durch
einen substantiellen mathematischen Durchbruch ausgelöst wird.

\medskip\noindent
\textbf{Serienstatus (v3.0, 26.05.2026; Kritikalitäts-Linsen-Update 31.05.2026).}
Fünfteilige Serie; alle drei bedingten Reduktionen unverändert seit v2.2.
Nur strukturelle Reorganisation; keine neuen Beweisbehauptungen. Eine Arbeitsphase
2026 ergänzt eine \emph{heuristische} Kritikalitäts-Linse
($\mathrm{RH}\Leftrightarrow\Lambda=0$), eine schärfere Lokalisierung der
gemeinsamen Wand und eine berechenbare Laguerre--P\'olya/Herglotz-Diagnostik
(Teil~IV~\cite{Geiger2026partIV}); diese gewichten die Prioritäten neu und schließen
vier \emph{konkrete Konstruktionen}, machen aber keine neue Beweisbehauptung.
\end{abstract}

{\footnotesize
\noindent\textbf{MSC 2020:} 11M26, 11M36, 47A75\\
\textbf{Schlüsselwörter:} Riemannsche Vermutung, bedingte Reduktion,
Atlas, Bewertung, offene Probleme, Ausblick
}

\clearpage


% =============================================================================
\section{Atlas-Übersicht}
\label{sec:atlas}
% =============================================================================

\begin{center}
\renewcommand{\arraystretch}{1.3}
\resizebox{\textwidth}{!}{%
\begin{tabular}{l|l|l|l|l}
\textbf{Route} & \textbf{Sektor} & \textbf{Status} & \textbf{Schlüsselergebnis} & \textbf{Paper} \\\hline
A (Even Dom.) & Even (Teil~II) & bedingt & A1--A8-Kette + GF1-Benchmark mit $C^* = 64\pi^2/3$ für $\lambda=5,7$; $\lambda=11$ offen & \cite{Geiger2026partII} \\
B (Hadamard) & Spektral (Teil~III) & bedingt & $A_\lambda \leq 1.0063$ unter Log-Krümmungskontrolle & \cite{Geiger2026partIII} \\
Z (CCM) & Spektral (Teil~III) & bedingt & C2ca.1 + C2ca.2 + C2ca.5; $g_0(m) = 2\pi/\log\gamma_m$ & \cite{Geiger2026partIII, Geiger2026zookeeper} \\
Routenübergreifend & Kombiniert (Teil~IV) & empirisch & gemeinsames Galerkin-Muster; 2026 Kritikalitäts-Linse + geschärfte gemeinsame Wand ($\Re s=1$-Fortsetzung) & \cite{Geiger2026partIV} \\
\end{tabular}}
\end{center}


% =============================================================================
\section{Ergebnis-Zusammenfassung}
\label{sec:results}
% =============================================================================

\subsection{Drei strukturelle Identifikationen}

\begin{enumerate}
  \item \textbf{Hadamard-Krümmungskonstante (RH-neutral):}
    \[
      c_\Xi = -\frac{1}{2}\sum_\rho
              \frac{1}{(\rho-\tfrac{1}{2})^2}
            \approx 0.0231.
    \]
    Unter der RH spezialisiert sich dies zu $\sum_{k\geq 1} 1/\gamma_k^2$.
    (Teil~III~\cite{Geiger2026partIII}.)

  \item \textbf{Randmoment-Konstante:}
    $C^* = 64\pi^2/3 \approx 210.55$,
    identifiziert als Kandidat über
    $16\int_0^\infty t^2/\sinh^2(t/2)\,dt$; der Faktor~$16$ wird
    als aus der DZT2-Leiternormalisierung stammend vermutet. Dies ist
    ein $\lambda=5,7$-Benchmark; $\lambda=3$ bleibt als eigener
    Niedrigband-Guardrail ausgewiesen, und $\lambda=11$ bleibt nach dem
    $N=200$-Punkt ein offener Diskriminator.
    (Teil~II~\cite{Geiger2026partII}.)

  \item \textbf{Hecke-normalisierte Lücke:}
    $g_0(m) \approx 2\pi/\log\gamma_m$, klassisch
    Montgomery--Odlyzko.
    (Teil~III~\cite{Geiger2026partIII}.)
\end{enumerate}

\subsection{Zwei bedingte Theoreme}

\begin{itemize}
  \item \textbf{Theorem B-1} (Teil~III): Unter uniformer
    Log-Krümmungskontrolle $H_\lambda(\delta) \leq K\delta^2$ mit
    $K \leq 1.8\times 10^{-3}$:
    \[
      A_\lambda(\delta)
      \leq B_\Xi \cdot \exp(K/4) \leq 1.0063,
    \]
    wobei $B_\Xi = \xi(0)/\xi(\tfrac{1}{2}) = 1.00579\ldots$
    analytisch ist.

  \item \textbf{Theorem GF1} (Teil~II): Unter eigenmodengewichteter
    Kernkonvergenz, die den effektiven Faktor~$16$ erzeugt:
    \[
      \mathrm{sLI}
      = \frac{64\pi^2}{3}\,\frac{\lambda}{NL} + o(\cdot).
    \]
    Dies ist für das nicht-ausnahmsartige konvergierte Regime formuliert;
    der $\lambda=3$-Sweep und der offene $\lambda=11$-Diskriminator
    verhindern eine All-$\lambda$-Lesart.
\end{itemize}


\subsection{Zwei isolierte offene Beweis-Slots}

\begin{enumerate}
  \item \textbf{B-1-Slot:} Uniforme Log-Krümmungskontrolle aus der
    Foundation-Lock-Struktur. Schließt Zweig~B.
  \item \textbf{GF1-Slot:} Kernkonvergenz + Faktor-16-Identifikation
    aus DZT2-Definitionen. Schließt den nicht-ausnahmsartigen
    Boundary-Moment-Benchmark von Zweig~A, mit separatem Niedrigband-Ast.
\end{enumerate}


% =============================================================================
\section{Lokalisierungstabellen}
\label{sec:localization}
% =============================================================================

\subsection{Li-Koeffizienten-Perspektive}

\begin{center}
\renewcommand{\arraystretch}{1.2}
\resizebox{\textwidth}{!}{%
\begin{tabular}{c|l|l|l}
\textbf{Schritt} & \textbf{Aussage} & \textbf{Status} & \textbf{Methode} \\\hline
S1 & Definition von $\xi(s)$, Funktionalgleichung & bewiesen & Definition \\
S2 & Li-Koeffizienten via Bombieri--Lagarias & bewiesen & \cite{BombieriLagarias1999} \\
S3 & Zerlegung $\lambda_n = \lambda_n^{(\infty)} + \lambda_n^{(\mathrm{fin})}$ & bewiesen & Teil~I \\
S4 & Archimedische Dominanz: $\lambda_n^{(\infty)} \sim \frac{n}{2}\log n$ & bewiesen & Stirling \\
\textbf{S5} & \textbf{$\lambda_n \geq 0$ für alle $n \geq 1$} & \textbf{OFFEN}${}^{\ast}$ & \textbf{$\Longleftrightarrow$ RH} \\
S6 & Alle Nullstellen auf $\Re(s) = 1/2$ & $\Leftrightarrow$ S5 & Li (1997) \\
S7 & Primzahlzählung: $\pi(x) = \mathrm{Li}(x) + O(\sqrt{x}\log x)$ & $\Leftarrow$ S6 & von Koch
\end{tabular}%
}
\end{center}

{\footnotesize\noindent${}^{\ast}$ S5 ist äquivalent zur RH (Li, 1997).
Die vorliegende Serie gibt nur bedingte Reduktionen.\par}


\subsection{Even-Dominance-Perspektive}

\begin{center}
\renewcommand{\arraystretch}{1.2}
\resizebox{\textwidth}{!}{%
\begin{tabular}{c|l|l|l}
\textbf{Schritt} & \textbf{Aussage} & \textbf{Status} & \textbf{Methode} \\\hline
E1 & Weil'sche explizite Formel & bewiesen & Weil (1952) \\
E2 & Weil'sche quadratische Form $\mathrm{QW}_\lambda$ & bewiesen & \cite{Connes2026} \\
E3 & Gerade/ungerade Zerlegung von $\mathrm{QW}_\lambda$ & bewiesen & Teil~II \\
E4 & Shift Parity Lemma: jede Primzahl bevorzugt gerade & bewiesen & Teil~II \\
E5 & Lückendiagnostik bei $33$ Werten, $\lambda \leq 1{,}300{,}000$ & Evidenz / bedingt & Teil~II \\
E5b & Euler--Maclaurin: $\rho^{\mathrm{EM}} < 0$ für $\lambda \geq 1.2\mathrm{M}$ & bewiesen & Teil~II \\
\textbf{E6} & \textbf{Kumulativer Schritt: Even Dominance für alle $\lambda$} & \textbf{bedingte Reduktion} & \textbf{= A6} \\
E7 & Connes' Theorem~6.1 + Hurwitz-Brücke $\Rightarrow$ RH & bedingt auf E6 & \cite{Connes2026}
\end{tabular}%
}
\end{center}


% =============================================================================
\section{Verbleibende offene Probleme}
\label{sec:open-problems}
% =============================================================================

\begin{remark}[OP1: Kumulative Even Dominance --- reduziert auf die untere Schranke des ungeraden Sektors (v2.2)]
Das kumulative Even-Dominance-Problem ist auf eine einzelne spektrale
Frage reduziert: eine quantitative untere Schranke für $\lambda_1^-(\lambda)$
(Asymptotic Variational Gap Conjecture, Teil~II).
\end{remark}

\begin{remark}[OP2 gelöst: Einfachheit von $\lambda_1^+$]
Für alle $33$ Zertifikatswerte durch Intervallarithmetik als einfach
zertifiziert. Die intra-gerade Lücke wächst wie $\sim\!\sqrt{\lambda}$
(Teil~II).
\end{remark}

\begin{conjecture}[OP3: Analytischer Beweis der Indexrigidität]
Der negative Trägheitsindex von $K_\sigma^{\mathrm{sym}}$ ist exakt~$2$
für alle hinreichend großen $N$ und
$\sigma \in (1, \sigma_{\max})$.
\end{conjecture}

\begin{conjecture}[OP4: Nuklearer Transferoperator für $\xi$]
Man konstruiere einen separablen Hilbertraum $\mathcal{V}$ und eine
nukleare Familie $s \mapsto \mathcal{L}_s$ mit
$\xi(s)/\xi(0) = \det(I - \mathcal{L}_s)$, Spektralradius $< 1$ für
$\Re(s) > 1/2$ und Eigenwert $1$ an den Nullstellen.
\end{conjecture}

\begin{conjecture}[OP5: Nevanlinna-Eigenschaft von $\log\det_2$]
Ist $\log\det_2(I - zR)$ genau dann eine Nevanlinna-Funktion, wenn alle
Nullstellen auf der reellen Achse liegen?
\end{conjecture}

\noindent OP2 ist gelöst. OP1 ist auf einen einzelnen asymptotischen
Schritt reduziert. OP3--OP5 bleiben strukturelle Fragen.


% =============================================================================
\section{Bewertung und Ausblick}
\label{sec:outlook}
% =============================================================================

\subsection{Was erreicht wurde}

\begin{enumerate}
  \item Eine \textbf{vollständige strukturelle Analyse} des
    eichtheoretischen Ansatzes (Teil~I: R1--R9, K1--K4).
  \item Das \textbf{Shift Parity Lemma} und \textbf{$33$ Finite-Range-Diagnostiken}
    mit Computerunterstützung (Teil~II).
  \item Die \textbf{bedingte Reduktion A1--A8} mit analytischem
    Rückgrat (M1$''$, Euler--Maclaurin, PNT-Transfer) und dem
    GF1-Boundary-Moment-Benchmark mit $C^* = 64\pi^2/3$ für
    $\lambda=5,7$ (Teil~II).
  \item Zwei \textbf{unabhängige spektrale bedingte Reduktionen}
    (Zweige~B und~Z, Teil~III).
  \item Die \textbf{routenübergreifende Synthese} mit Identifikation
    des gemeinsamen Galerkin-Musters und zweier isolierter Beweis-Slots
    (Teil~IV).
  \item \textbf{Fünfzehn ausgeschlossene Pfade} (spektral + CCM) und
    \textbf{fünf methodische Sackgassen} (Teile~III--IV).
  \item Eine 2026er \textbf{Kritikalitäts-Linse} ($\mathrm{RH}\Leftrightarrow\Lambda=0$,
    heuristisch), die die Routen zu exakt-strukturellen Mechanismen hin neu
    gewichtet; eine schärfere \textbf{Lokalisierung der gemeinsamen Wand} (die
    analytische Fortsetzung über $\Re s=1$); eine berechenbare
    \textbf{Laguerre--P\'olya/Herglotz-Diagnostik}; und vier Schließungen
    \emph{konkreter Konstruktionen} (keine Routen) --- alles ohne Änderung des
    Status der bedingten Reduktion (Teil~IV~\cite{Geiger2026partIV}).
\end{enumerate}


\subsection{Aktueller Atlas-Status}

Der Serientitel bleibt \glqq From Landscape to Atlas: Multi-Route
Cartography of an Ongoing Expedition Toward the Riemann Hypothesis.\grqq{}
Alle drei lebenden Routen sind bedingte Reduktionen; kein einzelner Zweig
hat unter dem Audit-Protokoll des Projekts das Niveau eines bedingten
Beweises erreicht.

\textbf{Heraufstufungs-Protokoll.} Die Titelmigration von \glqq Atlas\grqq{}
zu \glqq From Atlas to Proof\grqq{} wird durch ein Fünf-Tore-Audit
gesteuert:
(1)~einzelne benannte Hypothese $H$;
(2)~jeder Schritt von $H$ zur RH auf Theorem-Niveau;
(3)~zwei unabhängige Reviews einschließlich eines skeptischen Widerleger;
(4)~2--3 Tage Abkühlungsphase;
(5)~expliziter Decision-Record-Eintrag.
Standard: Atlas bleibt.


\subsection{Zukunftsszenarien}

\begin{itemize}
  \item \textbf{Ein Zweig stirbt} (seine offene Hypothese widerlegt):
    als tote Route dokumentiert; die anderen beiden laufen weiter.
  \item \textbf{Ein Zweig reift} (bedingungslos): dieser Beweis steht
    als RH; andere Zweige werden zur unabhängigen Verifikation beibehalten.
  \item \textbf{Beide/alle drei reifen}: selten, aber mathematisch
    wertvoll als Multi-Routen-Verifikation.
  \item \textbf{Keiner reift, aber partieller Fortschritt geht weiter}:
    das Forschungsprogramm wird fortgesetzt; die Landscape bewahrt das
    Audit. Dies ist der Standard-Ausblick.
\end{itemize}


\subsection{Ausblick}

Konkrete nächste Schritte (nicht zeitgebunden):
\begin{itemize}
  \item Rigorose Herleitung der zwei offenen Beweis-Slots (B-1
    Log-Krümmung, GF1 Faktor-16).
  \item Klärung des $\lambda=11$-GF1-Diskriminators, Validierung bei
    höheren $\lambda$ und theoremstarke Trennung des
    $\lambda=3$-Niedrigband-Asts.
  \item $\kappa_\infty$-Asymptotik für $\lambda \to \infty$.
  \item Fünf-Tore-Audit, falls ein Slot geschlossen wird.
  \item Entwicklung der exakt-strukturellen Linien (M~$=$~Bost--Connes\,$\leftrightarrow$\,Meyer;
    P-A~$=$~globale adelische Trace-Route), die die 2026er Kritikalitäts-Linse
    aufwertet, parallel zu den zwei Marge-artigen Beweis-Slots
    (Teil~IV~\cite{Geiger2026partIV}).
\end{itemize}


% =============================================================================
\section{Verbindungen zur breiteren Landschaft}
\label{sec:connections}
% =============================================================================

\subsection{Connes' Programm}

Die Arbeit in dieser Serie fügt sich direkt in Connes' spektralen Ansatz
zur RH~\cite{Connes2026} ein. Das Shift Parity Lemma (Teil~II) liefert
analytische Evidenz für die Even-Dominance-Hypothese (Theorem~6.1). Der
I-1 Normal-Family Reframe (Teil~III) zerlegt Wall~(ii) in eine einzelne
benannte Bedingung~(ii-a). Das CCM-Microcluster-Programm (Teil~III)
greift dieselbe Wand über das Fourier-Modell an.

\subsection{Nyman--Beurling--Baez-Duarte}

Das Li-Kriterium, das Weil'sche Positivitätskriterium, das
Nyman--Beurling-Kriterium und das Baez-Duarte-Kriterium sind alle
äquivalent zur RH. Direkte Implikationen zwischen ihnen (ohne den Umweg
über die RH) sind nur in zwei Fällen bekannt: Weil $\Rightarrow$ Li
(\cite{BombieriLagarias1999}) und Nyman--Beurling $\Rightarrow$
Baez-Duarte (trivial).

\subsection{Die Spurklassen-Obstruktion als universelles Muster}

Obstruktion~K4 (Spurklassen-Versagen bei $\Re(s) = 1/2$, Teil~I) tritt
quer durch QFT (UV-Divergenzen), statistische Mechanik (thermodynamischer
Limes) und NCG (Zetafunktions-Regularisierung) auf. Der
Even-Dominance-Ansatz umgeht dies, indem er mit endlich-dimensionalen
Galerkin-Blöcken statt mit unendlich-dimensionalen Transferoperatoren
arbeitet.


% =============================================================================
\section{Gesamtzusammenfassungs-Tabelle}
\label{sec:full-summary}
% =============================================================================

\begin{center}
\renewcommand{\arraystretch}{1.3}
\resizebox{\textwidth}{!}{%
\begin{tabular}{l|c|l}
\textbf{Beitrag} & \textbf{Status} & \textbf{Paper} \\\hline
Thermodynamischer Formalismus für $\zeta$ & etabliert & I \\
9 unbedingte strukturelle Resultate (R1--R9) & bewiesen & I \\
4 ausgeschlossene Pfade (K1--K4) & bewiesen & I \\
Shift Parity Lemma & bewiesen (analytisch) & II \\
Lückendiagnostik ($33$ Werte, $\lambda \leq 1{,}300{,}000$) & Evidenz / bedingt & II \\
Leading-Mode Cancellation ($c = 2+\sqrt{2}$) & bewiesen & II \\
Euler--Maclaurin: $\rho^{\mathrm{EM}}(13) < 0$ & bewiesen & II \\
GF1-Benchmark: $\mathrm{sLI} \sim C^*\lambda/(NL)$ für $\lambda=5,7$; $\lambda=3$ separat, $\lambda=11$ offen & bedingt (Thm GF1) & II \\
I-1 Normal-Family Reframe: (ii) $\to$ (ii-a) & rigoros & III \\
Foundation-Lock + konvergierte Empirie & rigoros + numerisch & III \\
Theorem B-1: $A_\lambda \leq 1.0063$ (bedingt) & bedingt & III \\
$c_\Xi$ Hadamard-Krümmung (RH-neutral) & analytisch & III \\
CCM drei C2ca-Inputs & bedingte Reduktion & III \\
6 ausgeschlossene spektrale + 9 ausgeschlossene CCM-Pfade & dokumentiert & III \\
Routenübergreifendes Galerkin-Muster & empirisch & IV \\
11 systematisch explorierte Alternativen & dokumentiert & IV \\
Hypothesenkorrektur-Kette (GF1 Iter 15--27) & dokumentiert & IV \\
\textbf{Kumulativer Schritt (A6)} & \textbf{bedingte Reduktion (v2.2)} & II \\
\textbf{Riemannsche Vermutung} & \textbf{bedingte Reduktion (v2.2)} & I--V
\end{tabular}%
}
\end{center}


\bigskip
\noindent\textbf{Danksagungen.} Berechnungen wurden auf einem Mac~Studio
(Apple M1~Max, 32~GB) und einem Hetzner CCX13 Cloud-Server (2~dedizierte
vCPU, 8~GB) durchgeführt. Die Intervallarithmetik-Bibliothek mpmath wurde
von Fredrik Johansson entwickelt.


\section*{KI-Offenlegung}

Dieses Manuskript wurde in intensiver Zusammenarbeit mit KI-Sprachmodellen
(Claude, Anthropic; Copilot, Microsoft/GitHub; Gemini, Google DeepMind)
erstellt. KI trug zu konzeptioneller Exploration, analytischer Arbeit,
Literaturrecherche, Schreiben und kritischer Überprüfung bei. Der Autor,
Lukas Geiger, konzipierte die Forschungsrichtung, brachte Fachexpertise
ein und übte die redaktionelle Letztverantwortung aus. Der Autor trägt die
alleinige wissenschaftliche Verantwortung für alle Inhalte, einschließlich
etwaiger Fehler.


\bibliographystyle{plain}
\bibliography{fst-rh-references}

\end{document}
